% !TEX root = ../thesis-main.tex

\chapter{Morse theory and metric algebraic geometry}
\label{C:metric-morse}

\begin{epigraph}
  When the blackbird flew out of sight, \\
  It marked the edge \\
  Of one of many circles.
\end{epigraph}

In the previous chapter, we presented a general theory for algebraically studying filtered spaces using the tools of persistent homology. In this section, we present tools from a slightly older theory, namely Morse theory, for studying filtered manifolds. We cover preliminaries on Morse theory and related studies of critical points of distance functions in metric algebraic geometry.

\section{A combinatorial encoding of topology}
\label{S:Morse-theory}

\begin{definition}
  Let $X$ be a smooth manifold and $f \colon X \to \R$ a smooth function. A \term[critical point]{critical point of $f$} is a point $x$ of $X$ where the differential $\partial_x f$ of $f$ is $0$. A \term{critical value} is the image $f(x)$ of a critical point.
  
  The \term[Hessian]{Hessian of $f$ at $x$} is the matrix of second partial derivatives of $f$ at $x$, denoted by $H(f)(x)$. A critical point $x$ of $X$ is \term[(non)degenerate critical point]{nondegenerate} if $H(f)(x)$ is nonsingular, and \textbf{degenerate} otherwise.
  
  Recall that the Hessian is a symmetric real matrix, and therefore is diagonalizable. The \term{index} of a nondegenerate critical point is the number of negative eigenvalues of its Hessian.
  
  A \term{Morse function} is a smooth function on a smooth manifold that has no degenerate critical points.
\end{definition}

The index of a nondegenerate critical point has the following interpretation in terms of normal forms: \ir{cite}

\begin{proposition}
  Let $X$ be a smooth manifold of dimension $n$, $f \colon X \to \R$ a smooth function, and $x \in X$ a nondegenerate critical point with index $\iota \in \{0, \ldots, n\}$. Then there exists a neighborhood $U$ of $x$ and a homeomorphism $\psi \colon \R^n \to U$ such that
  \begin{equation}
  \label{E:normal-form}
    f(\psi(t_1, \ldots, t_n)) = f(x) - t_1^2 - \cdots - t_\iota^2 + t_{\iota + 1}^2 + \cdots + t_n^2.
  \end{equation}
\end{proposition}

Viewing the filtration $f \colon X \to \R$ as a height function and considering the normal form \eqref{E:normal-form}, we can then consider the nondegenerate critical point $x$ to be a saddle point of $X$ where the manifold slopes downwards in $\iota$ directions and upwards in $n - \iota$ directions.

This algebraic encoding of the geometry of a filtered manifold can be developed further by the \term{Morse lemmas}, which we present as a theorem:

\begin{theorem}
  Let $X$ be a smooth manifold, $f \colon X \to \R$ a smooth function, and $s \leq t$ in $\R$. Recall the notation $X_{\leq t} = f^{-1}((-\infty, t])$. Then:
  \begin{enumerate}
    \item If the interval $[s, t]$ contains no critical values, then $X_{\leq s}$ is diffeomorphic to $X_{\leq t}$ and $X_{\leq t}$ deformation retracts onto $X_{\leq s}$.
    \item Suppose that there is a unique nondegenerate critical point $x \in X$ whose value $f(x)$ is in the interval $[s, t]$, and let $\iota$ be the index of $x$. Then $X_{\leq t}$ is homotopy equivalent to $X_{\leq s}$ with an $\iota$-cell attached.
  \end{enumerate}
\end{theorem}

\begin{example}
  \ir{classic donut}
\end{example}

In terms of homology, we have the following results \ir{cite}

\begin{definition}
  Let $X$ be a smooth manifold. For all integers $d \geq 0$, we call $b_d(X) \coloneq \dim H_d(X)$ the \term[Betti number]{$d^{\th}$ Betti number of $X$}.
\end{definition}

\begin{theorem}
  Let $X$ be a smooth manifold and $f \colon X \to \R$ a smooth function whose critical points are isolated and nondegenerate. For all integers $d \geq 0$, denote by $C^d$ the number of critical points of $f$ of index $d$. Then, for all $d \geq 0$, we have the inequality
  \[
    C^d - C^{d - 1} + \cdots + (-1)^d C^0 \geq b_d(X) - b_{d - 1}(X) + \cdots + (-1)^d b_0(X).
  \]
  In particular, for all $d \geq 0$, we get
  \[
    C^d \geq b_d(X).
  \]
\end{theorem}

When the set of critical points is easy to compute, this gives powerful upper bounds on the homology of the manifold $X$.

\section{Morse theory for continuous selections}

\begin{remark}
  Morse theory was originally developed for smooth function on smooth manifolds. However, it has been generalized in many directions, including Morse-Bott theory \ir{cite}, stratified Morse theory \ir{cite}, and discrete Morse theory \ir{cite}. In \ir{Paper D} we develop Morse theory for continuous selections, as presented in \cite{APS1997}.
\end{remark}

\ir{I'm starting to think that anything that is in the papers should not be written out here... UPDATE: only the new contributions, so continuous selections should stay here.}

\section{Distance functions between algebraic varieties}

For algebraic spaces, the notion of critical points of distance functions has been studied extensively in the past few decades, as part of the growing field of \term{metric algebraic geometry}.

\begin{definition}
  An \term{real algebraic variety} is the set of solutions of a system of multivariate polynomial equations over the real numbers. \ir{copied from Wikipedia, I guess...}
  
  Consider the space $\R^n$ for some $n \geq 1$ and an integer $d \geq 0$. Denote by $\cc{P}_{n, d}$ the space of real polynomials in $n$ variables, of degree $d$. Given polynomials $\vec{p} = p_1, \ldots, p_k$ in $\cc{P}_{n, d}$, we denote by $V(\vec{p})$ the algebraic variety that is the set of solutions of the equations $p_i(x) = 0$, for all $i \in \{1, \ldots, k\}$.
  
  The \term[(co)dimension]{dimension} of a variety $X$ is its dimension around any nonsingular point, seen as a manifold. We denote it by $\dim X$. Its codimension is $\codim V \coloneq n - \dim X$ where $n$ is the dimension of the ambient real or complex space.
  
  An algebraic variety $X$ is a \term{complete intersection} if there exist $\codim V$ many polynomials $p_1, \ldots, p_{\codim X}$ such that $X = V(p_1, \ldots, p_{\codim X})$.
\end{definition}

Algebraic geometers are often interested in integer invariants for algebraic varieties, and metric algebraic geometers are no exception. We present two of these \textit{degrees}:

\begin{definition}
  Let $X$ be a real algebraic variety and $u$ an arbitrary point in $\R^n$. Consider the function
  \[
    \dist_u \colon \left\{ \begin{array}{ccc}
      X & \to & [0, \infty) \\
      x & \mapsto & \| u - x \|_2^2
    \end{array} \right.
  \]
  defined as the squared Euclidean distance. It turns out that the number of critical points of $\dist_u$ is constant for a dense open subset of choices of $\R^n$: the \term[Euclidean distance degree (EDD)]{Euclidean distance degree} (\textbf{EDD}) \textbf{of $X$} is then this constant number \cite{DHOST2016}. An interpretation of the EDD is as a measure of the complexity of finding the closest point on $X$ to a given point $u$ in $\R^n$.
  
  Now let $X$ be a real algebraic variety. The \term[bottleneck]{bottlenecks of $X$} are pairs $(x, y)$ of distinct points of $X$ such that the line passing through $x$ and $y$ is normal (orthogonal) to $X$ at both points. Equivalently, bottlenecks are the nontrivial critical points of the squared distance function on $X$, that is, the critical points of the function
  \[
    \dist\big|_X \colon \left\{ \begin{array}{ccc}
      X \times X & \to & [0, \infty) \\
      (x, y) & \mapsto & \| x - y \|_2^2,
    \end{array} \right.
  \]
  with the extra condition that $x \neq y$. The \term[bottleneck degree (BND)]{bottleneck degree} (\textbf{BND}) \textbf{of $X$} is the number of bottlenecks, under some regularity conditions on $X$ \cite{DREW2020}.
\end{definition}

\begin{remark}
  Both of these notions fit within the framework of critical points of the distance function from an algebraic variety, restricted to another algebraic variety. In \ir{Papers D and E}, we generalize this notion to generic varieties.
\end{remark}